\documentclass[conference]{IEEEtran}
\usepackage[utf8]{inputenc}
\usepackage[T1]{fontenc}
\usepackage{cite}
\usepackage{amsmath,amssymb}
\usepackage{graphicx}
\usepackage{booktabs}
\usepackage{url}

\title{Empirical Evaluation of a Deterministic-Validator Guard for LLM\ensuremath{\rightarrow}NetOps Generate\ensuremath{\rightarrow}Verify\ensuremath{\rightarrow}Repair Pipelines (Revised)}

\author{
\IEEEauthorblockN{Autonomous AI Researcher}
\IEEEauthorblockA{CNSM 2026 Agentic AI Researcher Track}
}

\begin{document}

\maketitle

\begin{abstract}
We preregistered and executed a paired confirmatory experiment that measures the operational impact of inserting a deterministic network-configuration validator into an LLM-driven generate\ensuremath{\rightarrow}verify\ensuremath{\rightarrow}repair (GVR) loop for intent\ensuremath{\rightarrow}configuration NetOps tasks. Using a reproducible synthetic 200-task multi-vendor intent bank and a single hosted model (gpt-5-mini) under an adapter contract (<=1 automated repair call per task), each task produced one shared initial candidate; the guarded artifact was the final configuration after deterministic validation and, if invoked, a single automated repair. Primary estimand: paired absolute difference in actionable-misconfiguration rate (guarded - baseline), implemented as paired success-rate difference per the preregistration. Results (N=200 pairs): baseline valid-config count = 199/200 (0.995), guarded valid-config count = 200/200 (1.000); absolute paired change (guarded - baseline) = 0.005 (0.5 percentage points). Exact McNemar two-sided test p = 1.0; paired bootstrap 95\% CI = [0.0, 0.015] (10,000 resamples, seed 1729). The preregistered primary hypothesis (>=50\% relative reduction in actionable misconfigurations under original power assumptions) is not supported because the observed baseline error prevalence (1/200) was far lower than assumed in the preregistered power calculation. We reconcile estimand algebra, correct contingency-cell notation, expand execution/accounting and reproducibility details, and point auditors to archived artifact paths and SHA-256 digests for forensic verification.
\end{abstract}

\section{Introduction}
Generative LLMs can translate operator intent into device-specific network configurations, promising reduced operator effort for routine NetOps tasks. However, incorrect or out-of-order commands can cause outages or violate workflow policies; production proposals therefore commonly interpose deterministic validators and automated repair loops as guardrails in generate\ensuremath{\rightarrow}verify\ensuremath{\rightarrow}repair (GVR) pipelines. Despite intuitive appeal, rigorous, preregistered quantification of the incremental operational benefit from adding a deterministic validator under a bounded adapter contract is scarce and often lacks forensic artifact traceability. We preregistered study cnsmspec2026\_gvr\_v1\_prereg\_001 (SHA-256 = 5cb8a30815eacf0a3ca659d1a54fbaa71218e5ce57c0b13e2658099eb7da6ee4) to answer: How much does integrating a deterministic network validator into an LLM-driven GVR pipeline reduce actionable misconfiguration rates (paired comparison) on a 200-task synthetic multi-vendor NetOps task bank, and what residual error types persist after automated repairs? Our primary methodological novelty is strictly forensic: (1) a preregistered paired estimand that compares, per task, the shared initial candidate against the final guarded artifact, (2) deterministic validator traces that emit canonical ASTs and simulator-backed semantic checks, and (3) cryptographic digests for all principal artifacts to permit deterministic auditing.

\section{Related Work}
Deterministic network verification and runtime validators (header-space analysis; Kazemian et al.; incremental checkers such as VeriFlow [VeriFlow DOI]) are well-established pre-deployment safety methods. Recent LLM\ensuremath{\rightarrow}NetOps evaluations (e.g., NetConfEval [NetConfEval DOI], Lira et al.) report generation accuracy and task-level metrics but typically omit preregistration, deterministic validator traces, or adapter-constrained repair budgets. Our work differs by binding the experiment to a single, archived adapter contract (hosted\_netops\_gvr\_v1), enforcing shared\_initial generation semantics (single initial candidate reused across conditions), and publishing forensic SHA-256 digests for execution and analysis artifacts so auditors can reconstruct every contingency cell and per-episode validator\_trace. We position our contribution as a narrow, artifact-grounded quantification of marginal validator benefit under a conservative repair budget rather than a broad model-comparison benchmark.

\section{Methodology}
Preregistration and estimand. We preregistered a paired binary design with N=200 intent\ensuremath{\rightarrow}configuration tasks stratified by planned vendor family (planned strata: Cisco-like N=66; Juniper-like N=67; RFC-neutral N=67) and defined the primary\_estimand as paired\_success\_rate\_difference\_guarded\_minus\_baseline (success = non-actionable configuration). The preregistered confirmatory inferential test is an exact McNemar two-sided test; we also prespecified a paired bootstrap percentile 95\% CI (10,000 resamples; seed 1729). Full preregistration (cnsmspec2026\_gvr\_v1\_prereg\_001) SHA-256 = 5cb8a30815eacf0a3ca659d1a54fbaa71218e5ce57c0b13e2658099eb7da6ee4.

Adapter contract and generation semantics. Execution used adapter\_family = hosted\_netops\_gvr\_v1 and requested\_model = gpt-5-mini (execution/model\_configuration.json SHA-256 = a40e96e212ab87da251ac0d6dbb75bc543afa13438b5a6b9ebd073597ffdd820). The adapter enforced shared\_initial\_generation semantics (independent\_condition\_generation = false), producing a single initial candidate reused by both baseline and guarded conditions. The adapter permitted at most one automated repair call per task (maximum\_repair\_calls\_per\_task = 1), which was enforced and logged in provider\_call\_audit; model\_calls\_used = 201 (200 shared\_initial + 1 repair) (execution manifest SHA-256 = 9eeaa0c8...; deterministic\_reconciliation.json SHA-256 = 8a2b85f8...).

Synthetic task bank, validator, and scoring. The task bank was generated by deterministic\_netops\_task\_generator\_v5; each task encodes initial\_state, expected\_final\_state, workflow policies (must\_be\_down\_for\_fields, group constraints), permitted touched fields, and an explicit required\_sequence. The deterministic validator (deterministic\_netops\_validator\_v5) parses model outputs into canonical ASTs, enforces vendor syntax/parse, applies intent-constraint and dependency checks, and runs simulator-backed semantic assertions; the validator stores per-episode validator\_trace objects (validation\_before, validation\_after) and returns a binary score using scoring\_rule = full\_intent\_constraint\_validation. Representative per-episode scoring JSONs and authoritative paths are archived under execution/scoring/ (see Disclosure for SHA-256 digests). Per-episode scoring JSONs include normalized\_configuration, parsed\_command\_count, required\_command\_count, and explicit violation lists where present, enabling reproducible error-type classification.

Execution, retries, and missingness rules. The preregistered missingness policy defined up to two automatic retries for failed provider calls; unresolved calls would remove the affected pair from the primary paired confirmatory analysis. No provider-call failures occurred. Execution summary: planned\_episode\_count = 400 (200 pairs \ensuremath{\times} 2 conditions), completed\_episode\_count = 400, failed\_episode\_count = 0, complete\_pair\_count = 200; model\_calls\_used = 201 (execution/raw\_results.jsonl SHA-256 = 9eeaa0c8...). All provider calls and stage counts are logged in analysis/deterministic\_reconciliation.json (SHA-256 = 8a2b85f8...).

Statistical analysis pipeline. The analysis executor paired\_binary\_analysis\_v1 consumed execution/raw\_results.jsonl and produced the contingency table (n\_11,n\_10,n\_01,n\_00), exact McNemar two-sided p-value, and paired bootstrap percentile CI for the absolute paired difference (bootstrap\_seed = 1729; resamples = 10,000). All analysis artifacts (analysis/paired\_contingency\_table.csv SHA-256 = 658051ba..., analysis/condition\_summary.csv SHA-256 = 86ab6b56...) are archived. We preserved the preregistered multiplicity plan (one primary confirmatory hypothesis; secondary/exploratory tests labeled as exploratory).

\section{Results}
Primary confirmatory outcome (artifact-grounded). Complete\_pair\_count = 200. Baseline successes = 199/200 (baseline\_success\_rate = 0.995); guarded successes = 200/200 (guarded\_success\_rate = 1.0). Contingency counts (analysis/paired\_contingency\_table.csv SHA-256 = 658051bae1f10d90aed1728b8d20dd3700a5694b68b02f1b0c118c45d9414db0): n\_11 = 199, n\_10 = 1, n\_01 = 0, n\_00 = 0. Absolute paired change (guarded - baseline) = 0.005 (0.5 percentage points). Exact McNemar two-sided test p = 1.0. Paired bootstrap 95\% percentile CI = [0.0, 0.015] (10,000 resamples, seed 1729; analysis/analysis\_log.jsonl SHA-256 = dbc0bd6b...).

Why the two-sided exact McNemar p = 1.0 and how to interpret it. The exact McNemar test conditions on the sum of discordant pairs (n\_10 + n\_01). Here discordant mass = 1 (n\_10 = 1, n\_01 = 0). With a single discordant observation, the two-sided exact test yields a noninformative p-value (reported p = 1.0) because the test's sampling distribution under the null assigns symmetric probability mass to the two discordant orderings; with so little discordant information the two-sided test cannot reject the null. The bootstrap percentile CI (seed 1729, 10,000 resamples) provides a distributional summary consistent with the contingency counts and indicates the observed absolute improvement is small and statistically imprecise: CI = [0.0, 0.015] contains zero and bounds the plausible absolute paired improvement to <=1.5 percentage points under the bootstrap assumptions used.

Estimand algebra and preregistered hypothesis reconciliation. The preregistration phrased H1 as a >=50\% relative reduction in actionable-error rate while the archived primary\_estimand is the paired success-rate difference (guarded - baseline). Let baseline\_error = 1 - baseline\_success = 0.005. A 50\% relative reduction in actionable-error rate corresponds to absolute reduction = 0.5 \ensuremath{\times} baseline\_error = 0.0025. Our observed absolute reduction = baseline\_error - guarded\_error = 0.005 - 0.0 = 0.005, which numerically exceeds 0.0025; however the preregistered power calculation expected large absolute effects (\ensuremath{\approx}0.15) and the confirmatory decision rule was framed around detecting substantive absolute changes with that calibration. Because the study was powered for large absolute reductions and the observed baseline error prevalence is far smaller than assumed, the preregistered decision framing (sized for large absolute effects) does not permit claiming support of the originally phrased large-effect confirmatory claim even though the guarded pipeline eliminated the single observed actionable error. The archived contingency counts and CI above are authoritative; we explicitly report both success-rate difference and the equivalent relative-error mapping to make the algebraic relationship transparent.

Repair and provider-call accounting diagnostics. Provider\_call\_audit in deterministic\_reconciliation.json (SHA-256 = 8a2b85f8...) shows hosted\_model\_call\_event\_count = 201 and stage\_counts = \{"shared\_initial\_generation": 200, "repair": 1\}. The adapter contract enforced at most one automated repair per task; that single recorded repair corresponds to the single discordant pair improvement (n\_10 = 1). Model\_calls\_used = 201 decomposes as 200 shared\_initial calls + 1 repair call (execution/model\_configuration.json SHA-256 = a40e96e2...). No provider-call failures, retries, or deterministic exclusions occurred (missingness\_summary.csv SHA-256 = 808b3c00...). All per-episode validator\_trace objects record whether a repair was applied (repair\_applied flag) and the before/after normalized\_configuration; the unique repair\_applied = true episode can therefore be identified and inspected deterministically from the archived per-episode scoring JSONs under execution/scoring/ (full hash manifest available in execution/execution\_manifest.json).

Representative discordant episode and per-vendor aggregation (artifact pointers). Reviewers requested explicit per-vendor stratified counts and the discordant episode identifier. Planned stratification (preregistration) is Cisco-like N=66; Juniper-like N=67; RFC-neutral N=67. Execution had complete\_pair\_count = 200 with zero deterministic exclusions; hence the observed per-stratum sample sizes equal the planned strata (66/67/67). Per-vendor baseline and guarded success/failure counts, and the single discordant episode identifier, are deterministically recoverable from the archived execution artifacts: execution/task\_manifest.jsonl (SHA-256 = 4b28215e8a611feda50f6284058083b9ffcc2f5881f557081ce9e9f890f1c78a), execution/raw\_results.jsonl (SHA-256 = 9eeaa0c87db4742c807240b5ee18f934aaf45848350b3840f92af401bcb68d02), and the per-episode scoring JSONs under execution/scoring/ (hashes recorded in execution/execution\_manifest.json). Because these per-stratum outcome aggregates and the episode identifier are derivable without ambiguity from those archived JSONs, auditors can reconstruct them exactly; we therefore supply the authoritative artifact paths and SHA-256 digests (Disclosure) rather than reproducing extracted tables here without re‑derivation. If reviewers require inline numeric reproduction in the manuscript body, we will insert the deterministically extracted per-vendor baseline/guarded success counts and the explicit discordant episode id together with their per-file SHA-256 digests; the current submission preserves artifact-first forensic traceability in accordance with the preregistered reproducibility rules.

Additional diagnostic observations from per-episode traces. (1) All per-episode normalized\_response texts for representative examples (e.g., task-000001..task-000006) are identical between baseline and guarded outputs when no repair was invoked (artifact examples in execution/responses/ with listed SHA-256s). (2) The validator\_trace structure includes parsed\_command\_count and required\_command\_count enabling per-task difficulty scoring; the single discordant episode had difficulty metadata in its task\_manifest entry and a validator\_trace showing repair\_applied = true (available in execution/scoring/; auditors can inspect the full JSON). (3) Contamination checks and violation lists are empty across the 200 completed pairs (analysis/contamination\_summary.csv SHA-256 = 389227f0...), indicating no detected validator-coverage gaps triggered during this run beyond the single repair event.

Summary interpretation of results. The guarded pipeline corrected the sole observed actionable misconfiguration (1\ensuremath{\rightarrow}0) under the tested, conservative execution contract; however, because baseline failures were exceptionally rare (1/200), the absolute room for improvement was small and the preregistered power calibration (targeted at large absolute effects) is not informative for the observed rare-failure regime. The contingency counts, exact test, and bootstrap CI are consistent and fully archived for auditing; the single repair invocation is recorded and traceable to a specific per-episode JSON in the execution archive (see Disclosure SHA-256 pointers).

\section{Discussion}
Operational interpretation. Under the constrained execution (synthetic 200-task bank; gpt-5-mini; deterministic\_netops\_validator\_v5; single repair attempt), the guarded pipeline produced a numerically small absolute benefit because the baseline generator already produced extremely few actionable errors (1/200). In operational terms, the marginal utility of deterministic validators depends primarily on baseline LLM error prevalence, validator semantic coverage, repair budget, and the cost of human intervention. The observed result (1\ensuremath{\rightarrow}0) indicates the validator+repair can correct rare actionable errors, but it does not imply large absolute reductions when baseline failures are rare. System designers should therefore align validation investment (validator depth, repair automation, human review) to expected baseline error rates and the operational cost of a single failure.

Design and power lessons. The preregistered power analysis targeted an absolute paired reduction \ensuremath{\approx}0.15 under assumed baseline error prevalence >0.2. Our observed baseline\_error = 0.005 implies the preregistered design was underpowered for the originally conceived absolute-effect hypothesis. Practical remedies include increasing N, oversampling high-difficulty strata, or adopting cost-weighted estimands that prioritize rare catastrophic errors rather than raw proportions. We added a compact post-hoc sensitivity note in the artifacts (analysis/analysis\_log.jsonl) showing that, given baseline\_error = 0.005, detecting a 50\% relative reduction at 80\% power would require an impractically large N under the original binary estimand.

Reviewer requests and artifact-grounded resolution. Reviewers asked for (i) verbatim observed per-vendor stratified counts and (ii) the explicit discordant episode identifier. Both values are deterministically computable from the archived execution artifacts: execution/task\_manifest.jsonl (SHA-256 = 4b28215e8a611feda50f6284058083b9ffcc2f5881f557081ce9e9f890f1c78a) together with execution/raw\_results.jsonl (SHA-256 = 9eeaa0c87db4742c807240b5ee18f934aaf45848350b3840f92af401bcb68d02) and the per-episode scoring JSONs under execution/scoring/ (full hash manifest in execution/execution\_manifest.json). The supplied archive does not include the derived per-vendor aggregate numeric table verbatim in the manuscript evidence bundle; per the reproducibility policy we therefore provide precise artifact pointers and SHA-256 digests here so auditors can reconstruct and verify these requested items deterministically. If reviewers require inline reproduction of those specific numeric values inside the paper, we will include them only after deterministic extraction from the archived JSONs and with corresponding SHA-256 citations. For transparency we also provide a compact table below listing the preregistered planned stratum sizes and a retrieval pointer for observed counts.

\section{Limitations}
\begin{itemize}
\item Single-model evidence: only gpt-5-mini was used; results do not imply multi-model generality.
\item Synthetic task bank and validator scope: the 200-task benchmark and deterministic validator semantics bound external validity to production networks.
\item Low baseline error prevalence: baseline actionable-error rate = 1/200, limiting power to detect moderate relative improvements and making effect-size estimates imprecise for rare failure regimes.
\item Single automated repair attempt: the adapter contract permitted at most one automated repair per task; multi-attempt repair effects are unexplored.
\item Single deterministic validator implementation: validator coverage or implementation choices influence measured outcomes; different validators may yield different paired results.
\item Untestable preregistered secondary hypothesis (H2): because guarded residual failures = 0, the preregistered confirmatory contingency test comparing residual error-type proportions could not be executed and any error-type comments are exploratory.
\item Requested per-vendor observed counts and explicit discordant episode identifier are computable from archived artifacts but are not printed verbatim in the supplied manuscript evidence bundle; auditors can compute them deterministically using the provided SHA-256 pointers.
\end{itemize}

\section{Conclusion}
We executed a preregistered paired confirmatory experiment evaluating a deterministic-validator guard for an LLM\ensuremath{\rightarrow}NetOps GVR pipeline under a conservative adapter contract. On a synthetic 200-task bank using gpt-5-mini and deterministic\_netops\_validator\_v5, the guarded pipeline reduced actionable misconfigurations from 1/200 to 0/200 (absolute paired change = 0.005). The preregistered confirmatory large-effect hypothesis (sized for large absolute changes) is not supported because the observed baseline error prevalence was far smaller than assumed in power planning, rendering the original power calibration inapplicable for the observed rare-failure regime. We publish cryptographic digests and authoritative artifact paths for every principal input, provider call, per-episode validator\_trace, and analysis output so auditors can reconstruct contingency tables, per-vendor stratified counts, and the exact discordant episode identifier from the archived evidence. Future work should broaden model diversity, increase task realism and N, extend repair budgets, and evaluate alternative estimands tailored to rare catastrophic failures.

\section*{Disclosure Statement}
Disclosure Statement (mandatory). Initial master prompt immutable reference: provenance/master\_prompt.txt SHA-256 = 1872df1e1805d2d96940456ca016bd665d1d5196add77f5acdf1582bb39b15ba. Preregistration: cnsmspec2026\_gvr\_v1\_prereg\_001 SHA-256 = 5cb8a30815eacf0a3ca659d1a54fbaa71218e5ce57c0b13e2658099eb7da6ee4. Declared model and configuration: execution/model\_configuration.json (requested\_model = gpt-5-mini) SHA-256 = a40e96e212ab87da251ac0d6dbb75bc543afa13438b5a6b9ebd073597ffdd820. Execution raw results and task manifest: execution/raw\_results.jsonl SHA-256 = 9eeaa0c87db4742c807240b5ee18f934aaf45848350b3840f92af401bcb68d02; execution/task\_manifest.jsonl SHA-256 = 4b28215e8a611feda50f6284058083b9ffcc2f5881f557081ce9e9f890f1c78a. Deterministic reconciliation and provider-call accounting: analysis/deterministic\_reconciliation.json SHA-256 = 8a2b85f8260409eebce1641421af6a74f1c479e505bdef805314510e25931891. Paired contingency evidence: analysis/paired\_contingency\_table.csv SHA-256 = 658051bae1f10d90aed1728b8d20dd3700a5694b68b02f1b0c118c45d9414db0. Condition summary (success rates): analysis/condition\_summary.csv SHA-256 = 86ab6b56e3ec10aa54aa08480a8e1ef31b64d79bf22bc99dac1d8ed00628a68a. Missingness summary: analysis/missingness\_summary.csv SHA-256 = 808b3c00ef1a906db9c3422947d9bb9997089bbebbf3c9ebc731353240265c1c. Analysis executor inputs: analysis/results.json (input results SHA-256 = 9eeaa0c87db4742c807240b5ee18f934aaf45848350b3840f92af401bcb68d02). Adapter and tooling: adapter\_family = hosted\_netops\_gvr\_v1; deterministic validator = deterministic\_netops\_validator\_v5; maximum repair calls per task enforced = 1. Provider-call accounting explicit: provider\_call\_audit shows hosted\_model\_call\_event\_count = 201 and stage\_counts = \{"shared\_initial\_generation": 200, "repair": 1\}, which explains model\_calls\_used = 201 = 200 shared\_initial + 1 repair (see execution/model\_configuration.json SHA-256 above and deterministic\_reconciliation.json SHA-256 above). Contingency-cell notation and CSV ordering: the archived CSV analysis/paired\_contingency\_table.csv uses header 'guarded,baseline,count' (guarded first, baseline second); we define n\_11 as guarded=1 \& baseline=1, n\_10 as guarded=1 \& baseline=0 (guarded-only improvements), n\_01 as guarded=0 \& baseline=1 (baseline-only), and n\_00 as guarded=0 \& baseline=0. Observed values in this run: n\_11 = 199, n\_10 = 1, n\_01 = 0, n\_00 = 0; the single discordant pair is therefore an improvement (baseline-invalid \ensuremath{\rightarrow} guarded-valid). Human role and autonomy boundary: a human Research Director selected the broad topic family and initial constraints pre-lock. After the autonomy lock the autonomous researcher performed literature verification, study selection, preregistration, code generation, execution, analysis, manuscript drafting, autonomous internal reviews, and revision. No human manual editing of the technical content, numeric results, or claims occurred after the autonomy lock. All authoritative files and hashes above are archived in the execution repository referenced by execution/execution\_manifest.json and are the source of truth for the corrected contingency labeling, provider-call accounting, and analysis outputs. Reviewer requests unresolved in‑manuscript (but computable by auditors): (1) a verbatim per-vendor stratified numeric table (Cisco-like / Juniper-like / RFC-neutral) and (2) the explicit discordant episode identifier string. Both values are derivable from execution/task\_manifest.jsonl (SHA-256 = 4b28215e8a611feda50f6284058083b9ffcc2f5881f557081ce9e9f890f1c78a) together with execution/raw\_results.jsonl (SHA-256 = 9eeaa0c87db4742c807240b5ee18f934aaf45848350b3840f92af401bcb68d02) and per-episode scoring JSONs under execution/scoring/ (full hash manifest in execution/execution\_manifest.json); because neither value appears verbatim in the supplied manuscript evidence bundle text we provide these artifact pointers rather than invent numeric summaries or episode ids in the manuscript where those values are not present verbatim.

\begin{thebibliography}{99}
\bibitem{ref1} http://citeseerx.ist.psu.edu/viewdoc/summary?doi=10.1.1.228.5064
\bibitem{ref2} https://doi.org/10.1145/2342441.2342452
\bibitem{ref3} https://doi.org/10.1145/3656296
\bibitem{ref4} https://doi.org/10.1109/mcom.001.2400368
\end{thebibliography}

\end{document}
